\documentclass[11pt]{article}

\usepackage[preprint]{acl}
\usepackage{times}
\usepackage{latexsym}
\usepackage[T1]{fontenc}
\usepackage[utf8]{inputenc}
\usepackage{microtype}
\usepackage{inconsolata}
\usepackage{graphicx}
\usepackage{booktabs}
\usepackage{amsmath}
\usepackage[table]{xcolor}
\usepackage{multirow}
\usepackage{enumitem}

\setlength\titlebox{5cm}

% Anonymized for double-blind review
\title{A Scaled-Up Empirical Study of Syntactic Language Models}

\author{Pengcheng Wang\footnotemark[3]\footnotemark[4] \quad Zhexuan Zhang\footnotemark[3]\footnotemark[4] \quad Chao Lou\footnotemark[5] \quad Kewei Tu\footnotemark[3]\footnotemark[4]~\thanks{~~Corresponding author.} \\
\footnotemark[3]School of Information Science and Technology, ShanghaiTech University \\
\footnotemark[4]Shanghai Engineering Research Center of Intelligent Vision and Imaging\\
 \tt \{wangpch2024, zhangzhx2023, tukw\}@shanghaitech.edu.cn \\
\footnotemark[5]independent researcher \\
 \tt louchaobrn@outlook.com \\}

\begin{document}
\maketitle

\begin{abstract}
Syntactic language models (SLMs) that generate sentences along with their syntactic tree structures have shown promise in both syntactic generalization and downstream tasks. %%Some of these models enforce hierarchical attention patterns using structural attention masks. But two questions remain: whether the benefit comes from tree-structured input, structural attention masks, or both; and whether the approach scales to larger models and diverse tasks.
Existing studies, however, are relatively limited in scale regarding model size, training data size, and task diversity. In addition, the design space of tree linearization and attention masks has not been fully explored. In this paper, we conduct a scaled-up empirical study of 12 SLM variants, comparing different types of linearization and masks. 
%%We systematically compare 12 variants spanning different types of structured input and attention masks.
The experimental results show that SLMs with standard causal attention match or exceed the performance of those with structural masks across downstream and syntactic benchmarks, while preserving FlashAttention compatibility with zero architectural overhead. 
%The benefit comes from the input format, not the mask.
% We scale our method up with larger parameter sizes and diverse corpora, evaluating on more downstream tasks. 
% At the 1B scale, our method achieves the highest 11-task average.
We also observe the advantage of SLMs when scaling up at the 1B-parameter scale across 11 downstream tasks. Our code is available at \url{https://github.com/qlwpc/TG-Interpolation}.
\end{abstract}

\section{Introduction}

Standard causal language models process text as flat token sequences. While they can learn distributional proxies for syntax~\cite{hu-etal-2020-systematic}, hierarchical constituency structures---phrases composing into sentences---are left implicit. %%Providing explicit syntactic structure during pretraining could improve syntactic generalization and downstream reasoning. \cite{sartran-etal-2022-transformer}
Syntactic language models (SLMs)~\cite{sartran-etal-2022-transformer, hu-etal-2024-generative, murty-etal-2023-pushdown, zhao-etal-2024-dependency, huang2026GiLT} inject syntactic biases into Transformers~\cite{vaswani2017attention} using a simple technique: jointly modeling surface sentences and a linearized version of their syntactic parse trees. 
%%Prior work on SLMs has explored several directions.
%%Transformer Grammars~\cite{sartran-etal-2022-transformer} serialize constituency trees into token sequences and add structural attention masks, allowing closing non-terminals to compose their child constituents. They demonstrate syntactic gains at $\sim$100M scale, but don't explain whether the benefit comes from the tree-structured input tokens, the structural attention mask, or both.
They typically employ \emph{structural} attention masks that allow a non-terminal to compose its child tokens once they are generated by focusing its attention on them.
\citet{zhao-etal-2025-systematic} conduct a systematic empirical comparison of various design choices for SLMs, but they do not fully explore the complete design space of tree linearization and attention masks.
%%linearization and composition mechanisms. 
%%However, their comparison is restricted to composition strategies \textit{within} structural-mask models, and do not compare against baseline models with an equivalent amount of computation.
Moreover, all the previous studies of SLMs are limited in scale regarding model size (up to 300M), training data size (up to 9B tokens), and task diversity (mostly on 2-3 tasks).

In this paper, we conduct a more scaled-up empirical study of SLMs (model size up to 1B, training data up to 100B, and a more contemporary suite of 11 downstream tasks). In our study, we also investigate novel designs of tree linearization and attention masks, resulting in 12 SLM variants.
% In this paper, we address three unanswered questions: 
% \begin{enumerate}
%     \item \textbf{Disentanglement}: Does the benefit come from tree-structured input, the structural attention mask, or both? Prior work conflates them.
%     \item \textbf{Scalability and practicality}: Do SLMs remain beneficial at larger scales and on diverse corpora? % Can we drop the inefficient structural masks?
%     \item \textbf{Explicit vs.\ implicit}: How does explicit tree structure compare to implicit structure at matched training budgets?
% \end{enumerate}
% %
% \begin{description}
%     \item[Structural Attention] Tree input with STACK / COMPOSE attention masks (TG) or Nomask. This paradigm requires custom kernel to accelerate attention computing.
%     \item[Causal Attention] Tree input with standard causal attention.
%     \item[Hybrid Attention] Tree input with head-level mask mixing--- splitting half the attention heads using one structure, and the other half using a different structure. This paradigm also requires custom kernel to accelerate attention computing.
%     %\item[Latent-Structure:] We also address the hypothesis that the advantages of SLMs over vanilla LMs come from extra computation resulted from additional input tokens. Specifically, we compare SLMs with vanilla LMs that have dummy tokens inserted into the input sequence ~\cite{goyal2024think}.% Flat input interleaved with pause tokens and causal attention (Pause-1, Pause-2)~\cite{goyal2024think}.%
% \end{description}
%
%%We systematically compare \textbf{12 variants} spanning three attention mechanisms and several tree linearizations variants. 
Our findings are as follows: (i) SLMs with standard causal attention match or surpass structural-mask variants across downstream and syntactic benchmarks, indicating that the benefit mainly comes from modeling linearized syntactic trees rather than specialized masks; they also preserve native FlashAttention compatibility with zero architectural overhead. (ii) The advantage of SLMs persists when scaling to 1B parameters and more downstream tasks on diverse data, suggesting that explicit modeling linearized syntactic trees remains useful beyond small, narrow settings.

% TODO
In summary, our contributions include:
\begin{itemize}[noitemsep]
    \item We scale SLM evaluation along model size, training data, and task diversity, verifying their advantage at 1B parameters on 11 downstream tasks. %and introducing terminal-format evaluation for comparison.
    \item We systematically examine 12 controlled variants with varying tree linearization and mask strategies, disentangling their impact on model performance.
    \item We show that causal masks outperform structural masks while preserving FlashAttention compatibility, and provide possible explanation‌ of the gains through input-side syntactic supervision, learned CNT bottlenecks, and adaptive computation allocation.
\end{itemize}

\section{Related Work}

% \subsection{Syntactic Language Models}

Prior work on syntactic language models (SLMs) has explored a broad range of ways to incorporate linguistic structure into neural sequence modeling. Transformer Grammars (TG)~\cite{sartran-etal-2022-transformer} encode constituency trees as a sequence of actions and enforce hierarchical attention patterns. A following study extends the idea to dependency trees~\cite{zhao-etal-2024-dependency}. Building on this line of work, \citet{zhao-etal-2025-systematic} provide the most systematic comparison of SLM design choices including those of tree linearization and composition mechanisms. These studies show that syntactic inductive bias can improve language modeling and syntactic evaluation, but they typically couple modeling of linearized trees with specialized attention mechanisms.

Other work incorporates syntactic inductive bias through alternative architectural mechanisms rather than explicit representations. Examples include latent structure induction during pretraining~\cite{hu-etal-2024-generative}, stack-based recursive state tracking~\cite{murty-etal-2023-pushdown}, and dependency-graph-informed attention mechanisms~\cite{huang2026GiLT}.

% \subsection{Pause Tokens and Latent Computation}

% A separate line of work improves language models by giving them more computation. Pause tokens~\cite{goyal2024think} append learnable tokens to the input, delaying output prediction. 
% We include them as a strong baseline with comparable computation with SLMs. 
% Scaling Latent Reasoning~\cite{zhu2025scaling} extends pause-token logic to ``looped'' LMs where the same block iterates multiple times per token. 
% This positions pause tokens within a broader trend toward latent computation. 

% \subsection{Understanding Why Structure Helps}

% \citet{sartran-etal-2022-transformer} argue that structure helps by providing syntactic supervision: non-terminals constrain plausible next words, while TG-style restricted attention promotes composed subtree representations. Such structure improves syntactic generalization, but may trade off against lexical copying and long-context memory.

% % Psycholinguistic probes provide evidence that language models can encode syntactic state. \citet{futrell-etal-2019-neural} examine LMs as psycholinguistic subjects and find evidence for representations of syntactic state. More recent structural-priming evidence from multilingual language models further supports the view on LLMs~\cite{michaelov-etal-2023-structural}.

% This interpretation is also consistent with probability-based syntactic evaluation: \citet{hu-etal-2026-string} show that language-model string probability correlates with grammaticality judgments, supporting the use of log-probability as a syntactic signal. This motivates our use of document-level PPL and terminal-format scoring.

\section{Variants of Syntactic Language Models}

\subsection{Syntactic Language Models}

Let $\mathbf{x} = (x_1, \ldots, x_N)$ be a sentence of $N$ terminal tokens, and let $\mathbf{y}$ be a labeled constituency parse tree over $\mathbf{x}$, where each node carries a phrase category (e.g., \texttt{S}, \texttt{NP}, \texttt{VP}) and spans a contiguous subsequence of $\mathbf{x}$.
A syntactic language model defines a joint distribution $p(\mathbf{x}, \mathbf{y})$, factorized autoregressively over a linearized token sequence $\mathbf{z} = (z_1, \ldots, z_T)$:
\begin{equation}
p(\mathbf{z} := \textsc{Linearize}(\mathbf{x}, \mathbf{y})) = \prod_{t=1}^{T} p(z_t \mid z_{<t})
\label{eq:joint}
\end{equation}
where $\textsc{Linearize}$ maps the sentence--tree pair into a flat sequence. Under this formulation, prior SLMs can be understood as different ways of instantiating the same autoregressive modeling problem, varying in the \textbf{linearization} used to serialize the tree $\mathbf{y}$ into $\mathbf{z}$, and the \textbf{architecture} used to parameterize the conditional probabilities $p(z_t \mid z_{<t})$. In Transformer-based SLMs, this architectural choice is largely realized through the \textbf{attention mechanism}, which governs how syntactic relationships among tokens in 
$\mathbf{z}$ are processed.

% Crucially, while Eq.~\ref{eq:joint} defines the modeling objective, existing work differs in \textit{which tokens} the model predicts.
% When $\mathbf{z}$ contains structural tokens, the model may be trained to predict all tokens in $\mathbf{z}$ (including ONT and CNT).


%\paragraph{Latent-Structure} (flat input + learned ``thinking'' tokens). Rather than providing tree structure explicitly, the model receives a flat sequence $\mathbf{z} = (x_1, p_1, \ldots, p_K, x_2, p_1, \ldots, p_K, x_3, \ldots)$ where $\{p_k\}$ are learnable pause tokens~\cite{goyal2024think}.
%The model learns implicit structure through these extra positions with standard causal attention.
%Structure is learned rather than provided, at the cost of additional computation.


% Table~\ref{tab:paradigms} summarizes the four paradigms.

% \begin{table}[t]
%   \centering
%   \resizebox{\columnwidth}{!}{%
%   \begin{tabular}{lccc}
%     \toprule
%     Paradigm & Linearization & Attention & Fast Attn \\
%     \midrule
%     Mask-Structure & Tree-structured & Structural & No \\
%     Sequence-Structure & Tree-structured & Causal & Yes \\
%     Hybrid-Structure & Tree-structured & Hybrid & No \\
%     Latent-Structure & Flat + pause & Causal & Yes \\
%     \bottomrule
%   \end{tabular}%
%   }
%   \caption{Four paradigms as choices of linearization and attention.}
%   \label{tab:paradigms}
% \end{table}

\subsection{Tree Linearization}
\label{sec:tree_linearization}



\begin{table*}[t]
\centering
\renewcommand{\arraystretch}{1.2}
\resizebox{0.94\textwidth}{!}{%
\setlength{\tabcolsep}{4pt}
\begin{tabular}{lcccp{9cm}}
\toprule

\textbf{Name} & \textbf{Left($\mathbf{v}$)} & \textbf{Right($\mathbf{v}$)} & \textbf{Post-process} & \textbf{Example} \\
\midrule
\rowcolor{gray!15} 
\multicolumn{5}{l}{\textbf{Linearization}} \\
$\textsc{Lin}^1$ & ONT($v$) & CNT($v$) & - & \texttt{(S Cats (VP chase mice VP) S)} \\
$\textsc{Lin}^2$ & ONT($v$) & 2$\times$CNT($v$) & - & \texttt{(S Cats (VP chase mice VP) VP) S) S)}  \\
$\textsc{Lin}^3$ & ONT($v$) & 3$\times$CNT($v$) & - & \texttt{(S Cats (VP chase mice VP) VP) VP) S) S) S)}  \\
$\textsc{Lin}^1_{\text{--ont}}$ & $\emptyset$ & CNT($v$) & - & \texttt{Cats chase mice VP) S)}  \\
$\textsc{Lin}^1_{\text{--merge}}$ & ONT($v$) & CNT($v$) & Merge CNTs & \texttt{(S Cats (VP chase mice )} \\\midrule
\rowcolor{gray!15} 
\multicolumn{5}{l}{\textbf{Non-structural}} \\
$\textsc{Lin}^0$ & $\emptyset$ & $\emptyset$ & - & \texttt{Cats chase mice} \\
$\textsc{Lin}^1_{\text{shuf}}$  &  ONT($v$) & CNT($v$) & Shuffle & \texttt{(S VP) Cats (VP chase S) mice} \\
Pause-1 & $\emptyset$ & $\emptyset$ & Insert 1 <PAUSE> & \texttt{Cats <PAUSE> chase <PAUSE> mice} \\
Pause-2 & $\emptyset$ & $\emptyset$ & Insert 2 <PAUSE> & \texttt{Cats <PAUSE> <PAUSE> chase <PAUSE> <PAUSE> mice} \\
\bottomrule
\end{tabular}%
}
\caption{Overview of linearizations and non-structural token sequences used in this paper. ``Post-process'' refers to operations applied after linearization. ONT and CNT denote opening (e.g., \texttt{(S} and \texttt{(VP}) and closing non-terminal tokens (e.g., \texttt{S)} and \texttt{VP)}), respectively. $v$ indicates that the emitted token carries a node label. "Merged CNT" signifies that consecutive CNTs (ignoring labels) are collapsed into a single \texttt{)} token. Finally, $\textsc{Lin}^1_{\text{shuf}}$ is generated by shuffling the $\textsc{Lin}^1$ token sequence while preserving the original terminal order.}
\label{tab:linearizations}
\end{table*}


The function $\textsc{Linearize}(\mathbf{x}, \mathbf{y})$ embeds sentence $\mathbf{x}$ and tree $\mathbf{y}$ into token sequence $\mathbf{z}$. We explore several variants of linearization, all of which share a simple recursive form:
\begin{align}
    \textsc{Lin}(v) &:= \textsc{Left}(v) \notag \\
    &\quad \circ [\, \textsc{Lin}(c) \text{ for } c ~\textsc{in} ~\textsc{children}(v) \,] \notag \\ 
    &\quad \circ \textsc{Right}(v) \\[1ex]
    \textsc{Lin}(t) &:= t
\end{align}
where $v\in \mathbf{y}$ is a node, $t\in\mathbf{x}$ is a terminal token, and $\circ$ denotes concatenation. $\textsc{Left}(v)$ and $\textsc{Right}(v)$ produces \emph{structural tokens} to enclose the linearized children. $\textsc{children}(v)$ returns the ordered children of $v$ in the tree. 

Table~\ref{tab:linearizations} summarizes all linearization variants that we study. $\textsc{Lin}^1$, $\textsc{Lin}^2$ and $\textsc{Lin}^3$ are three variants of linearizations, differing only in the number of structural tokens emitted by $\textsc{Right}(v)$. Specifically, $\textsc{Lin}^1$ generates sequences resembling the widely used bracket notation for constituency trees, but with additionally labeled closing brackets. $\textsc{Lin}^2$ was proposed by \citet{sartran-etal-2022-transformer} to define their structural attention mask (referred to as STACK/COMPOSE in our paper). $\textsc{Lin}^3$ is a variant designed to test whether providing more closure tokens helps. We also include two additional variants, $\textsc{Lin}^1_{\text{--ont}}$ and $\textsc{Lin}^1_{\text{--merge}}$, to enable a more systematic comparison. %We select these variants to methodically probe the core design axes of syntax-aware linearization: structural token budget, opening tokens, and per-constituent closure. 
Other variations (e.g., different traversal orders) are not included because they would introduce confounds that distract from our central goal: disentangling tree linearization from the attention mask design. %characterizing emergent adaptive computation.

Different linearizations introduce varying numbers of structural tokens, inherently altering the total sequence length for training and inference (Table~\ref{tab:token-budgets}). %For instance, $\textsc{Lin}^1$ and $\textsc{Lin}^2$ produce $2.4\times$ and $3.2\times$ more tokens per raw document, respectively, compared to the baseline $\textsc{Lin}^0$. 
To separate performance gains due to syntactic structures from gains due simply to increased computation, we include three token-budget-matched baselines alongside the naive terminal sequence. First, $\textsc{Lin}^1_{\text{shuf}}$ shuffles the structural tokens to serve as a control for the necessity of valid structural information. The other two, PAUSE-1 and PAUSE-2, insert dummy $\text{<pause>}$ tokens~\cite{goyal2024think} uniformly throughout the sequence.


% The core sequence-structured linearization, which we call $textsc{Lin}^1$ (single CNT), produces one ONT and one CNT per constituent:
% \begin{multline}
% \textsc{Lin}^1(v) = \text{ONT}(\ell(v)) \; \circ \; \textsc{Lin}^1(\text{children}(v)) \\
% {}\circ\; \text{CNT}(\ell(v))
% \label{eq:lin1}
% \end{multline}

% where $\circ$ denotes concatenation. For the full tree rooted at $v_{\text{root}}$, $\mathbf{z} = \textsc{Lin}^1(v_{\text{root}})$.
% The standard TG approach~\cite{sartran-etal-2022-transformer} uses a doubled-closure variant:

% \begin{multline}
% \textsc{Lin}^2(v) = \text{ONT}(\ell(v)) \; \circ \; \textsc{Lin}^2(\text{children}(v)) \\
% {}\circ\; \text{CNT}_1(\ell(v)) \; \circ \; \text{CNT}_2(\ell(v))
% \label{eq:lin2}
% \end{multline}

% where $\text{CNT}_1$ and $\text{CNT}_2$ are identical token types (duplicated tokens).
% Under mask-structure attention, $\text{CNT}_1$ attends to constituent-internal tokens (stack) while $\text{CNT}_2$ attends to $\text{CNT}_1$ and prior closures (composition).
% Under causal attention ($\textsc{Lin}^2$ with Sequence-Structure), the model must learn to use this duplicated capacity.


\begin{table}[htbp]
  \centering
  % \small
  \setlength{\tabcolsep}{3.5pt}
  \begin{tabular}{lccc}
    \toprule
    Dataset & $\textsc{Lin}^0$ & $\textsc{Lin}^1$ & $\textsc{Lin}^2$ \\
    \midrule
    BBC News & $\sim$10B & $\sim$24.7B & $\sim$32B \\
    FineWeb-Edu 100BT & $\sim$98B & $\sim$233B & $\sim$301B \\
    \bottomrule
  \end{tabular}
  \caption{Training token counts for each linearization. Structural tokens inflate the token budget 2.4--3.2$\times$ per raw document.}
  \label{tab:token-budgets}
\end{table}



% \subsection{Linearization Variants}
% From these primitives, we define the full space of linearization variants studied in this work (Table~\ref{tab:variants}):

% \begin{description}
%     \item[$\textsc{Lin}^0$ (Terminal):] $\mathbf{z} = \mathbf{x}$, no structural tokens. The standard LM baseline.
%     \item[$\textsc{Lin}^1$ (Tree):] Eq.~\ref{eq:lin1}. One CNT per constituent.
%     \item[$\textsc{Lin}^2$ (TGTree):] Eq.~\ref{eq:lin2}. Two identical CNTs per constituent.
%     \item[$\textsc{Lin}^3$ (Tree-TripleCNT):] Three CNTs per constituent; tests whether more closure tokens help.
%     \item[$\textsc{Lin}^1_{\text{--ont}}$ (Tree-NoONT):] $\textsc{Lin}^1$ without ONT tokens, ablating opening non-terminals.
%     \item[$\textsc{Lin}^1_{\text{--merge}}$ (Tree-Compress):] Consecutive identical CNTs merged into one, ablating per-constituent closure.
%     \item[$\textsc{Lin}^1_{\text{shuf}}$ (Tree-Shuffle):] ONT and CNT labels randomly permuted, testing whether the model uses label identity vs.\ position.
% \end{description}

\begin{table*}[t]
  \centering
  % \small
  \setlength{\tabcolsep}{6pt}
  \begin{tabular}{cllll}
    \toprule
    \# & Variant & Linearization & Attention  \\
    \midrule
    1 & Terminal & $\textsc{Lin}^0$ & Causal  \\
    2 & Tree & $\textsc{Lin}^1$ & Causal \\
    3 & TGTree & $\textsc{Lin}^2$ & Causal  \\
    4 & TG & $\textsc{Lin}^2$ & STACK/COMPOSE  \\
    5 & TGNomask & $\textsc{Lin}^2$ & Nomask  \\
    6 & TGNomask-Aug$^*$ & $\textsc{Lin}^2$ & Nomask  \\
    7 & Tree-NoONT & $\textsc{Lin}^1_{\text{--ont}}$ & Causal \\
    8 & Tree-Compress & $\textsc{Lin}^1_{\text{--merge}}$ & Causal \\
    9 & Tree-TripleCNT & $\textsc{Lin}^3$ & Causal  \\
    10 & Tree-Shuffle & $\textsc{Lin}^1_{\text{shuf}}$ & Causal   \\
    11 & TGNomask-Mix-TG & $\textsc{Lin}^2$ & Mixing heads (Nomask + STACK/COMPOSE)  \\
    12 & TGTree-Mix-TG & $\textsc{Lin}^2$ & Mixing heads (Causal   + STACK/COMPOSE)  \\
    \bottomrule
  \end{tabular}
  \caption{Full specification of all 12 syntactic variants. ONT = opening non-terminal, CNT = closing non-terminal. ``Hybrid'' denotes mixing-head models. $^*$TGNomask-Aug differs from TGNomask: CNT$_2$ remains visible to subsequent attention in Aug, while TGNomask prevents subsequent tokens from attending to it.}
  \label{tab:variants}
\end{table*}

\subsection{Structural Attention}
\label{sec:structural_attn}

A Transformer-based SLM can also incorporate structural information directly into the model itself via specially-designed attention mechanisms. We study three attention mechanisms as listed below. 

%Structural mask variants (TG, TGNomask, TGNomask-Aug, ...) use $\textsc{Lin}^2$ as their linearization but differ in attention---STACK/COMPOSE (TG), \textbf{Nomask} (TGNomask, TGNomask-Aug), or head-level mask mixing.

\paragraph{Causal Attention} The joint $p(\mathbf{x}, \mathbf{y})$ is modeled through standard autoregressive factorization over $\mathbf{z} = \textsc{Linearize}(\mathbf{x}, \mathbf{y})$ with \textit{no} modification to the attention mechanism.
Constituency structures are encoded entirely in the token sequence; the model learns to route information through Opening Non-Terminal (ONT) and Closing Non-Terminal (CNT) using standard causal attention.
%This paradigm Zero architectural changes relative to standard LMs.

% \paragraph{Structural Attention} The structural attention mechanism factorizes $p(\mathbf{x}, \mathbf{y})$ autoregressively over $z = \textsc{Lin}^2(\mathbf{x}, \mathbf{y})$ and replaces standard causal attention with input-dependent structural masks.
% The original TG approach~\cite{sartran-etal-2022-transformer} uses STACK/COMPOSE attention. % a CNT$_1$ token aggregates all child tokens within the closing constituent via a stack operation; CNT$_2$ (a duplicate of CNT$_1$) implements recursive composition.
% \citet{zhao-etal-2025-systematic} systematically compare structural masks within this paradigm and find that SLMs without sub-constituent masks perform better. We follow their setting as \textbf{Nomask}.
\paragraph{Structural Attention}
Structural attention factorizes $p(\mathbf{x}, \mathbf{y})$ autoregressively over the linearized tree sequence $z = \textsc{Lin}^2(\mathbf{x}, \mathbf{y})$ and replaces standard causal attention with structure-dependent attention masks. In the original Transformer Grammar (TG) formulation~\cite{sartran-etal-2022-transformer}, each CNT is duplicated into two actions. The first action performs \textsc{COMPOSE} attention, which attends only to the sub-constituents inside the current constituent and computes a single composed representation. The second action performs \textsc{STACK} attention, which restricts attention to the active stack of previously completed constituents. After composition, the original sub-constituents are masked from future attention, forcing subsequent predictions to access them only through the composed representation. \citet{zhao-etal-2025-systematic} isolate this masking strategy as a separate design choice and show that removing the sub-constituent masking yields consistently better performance. Following their terminology, we adopt the \textbf{Nomask} setting, where composed sub-constituents remain directly accessible after composition.

\paragraph{Hybrid Attention} Hybrid attention mechanism factorizes $p(\mathbf{x}, \mathbf{y})$ autoregressively over $z = \textsc{Lin}^2(\mathbf{x}, \mathbf{y})$ and replaces standard causal attention with \textbf{head-level mask mixing}: rather than applying a single mask type across all attention heads, different heads receive different masks. The hybrid attention enables the model to combine multiple structural attention within a single Transformer layer.
Since enumerating all possible head‑wise mask assignments is computationally infeasible, we restrict ourselves to a straightforward setting: half of the attention heads receive one type of structural mask, and the remaining heads use another (either causal or structural). This configuration provides a tractable way to study hybrid attention.

Table~\ref{tab:variants} summarizes all the 12 variants across different designs of tree linearizations and attention masks studied in our framework.

\section{Terminal-Format Evaluation}
\label{sec:terminal-format-evaluation}

When applying a language model to multiple-choice (MC) downstream tasks, the standard protocol scores each answer option by summing the model's log-probability over all continuation tokens.
For SLMs, continuations contain structural tokens (ONT/CNT) interleaved with terminal tokens. The model typically assigns substantially higher probability to these structural tokens---on average $\sim$2.7 nats higher than terminal tokens, corresponding to $\sim$14$\times$ higher per-token probability (see Appendix~\ref{sec:decom} for a detailed decomposition).
%%This structural token preference dominates the aggregate MC score: summing log-probability over the full continuation mixes terminal content prediction with the model's intrinsic preference for bracket structure.
As a result, structural token probabilities dominate the aggregate MC score, compromising the evaluation of terminal content tokens.

We propose terminal-format evaluation which decomposes the MC score into its terminal and non-terminal components and uses only the terminal component for answer ranking.
Formally, for a continuation $\mathbf{c}$ produced by a syntactic model, let $\mathcal{T}(\mathbf{c})$ be the set of terminal indices and $\mathcal{N}(\mathbf{c})$ the set of non-terminal indices.
The \textit{full score} and \textit{terminal score} are:
\begin{align}
S_{\text{full}}(\mathbf{c}) &= \sum_{t \in \mathcal{T} \cup \mathcal{N}} \log p(c_t \mid c_{<t}) \\
S_{\text{term}}(\mathbf{c}) &= \sum_{t \in \mathcal{T}} \log p(c_t \mid c_{<t})
\label{eq:decom}
\end{align}
MC accuracy is computed by selecting the answer with the highest score; terminal-format evaluation uses $S_{\text{term}}$ instead of $S_{\text{full}}$.
% For terminal baselines ($\textsc{Lin}^0$), $\mathcal{N}(\mathbf{c})$ is empty and the two scores are identical---the decomposition has no effect.
% For pause-token models, pause positions are similarly excluded from $S_{\text{term}}$.

We report evaluating both full scores and terminal scores results for all SLMs.
Appendix~\ref{sec:decom} presents a full decomposition analysis across eight MC tasks, validating that terminal-format evaluation is a more effective method.


\section{Experiments}

\subsection{Setup}

We train all models from scratch for one epoch using the same OLMo-based~\cite{groeneveld-etal-2024-olmo} decoder-only Transformer family across scales. Learning rates and batch sizes follow the Step Law~\cite{li2025predictable}. Causal attention models use FlashAttention-2~\cite{dao2023flashattention2}, and structural-mask models require FlexAttention~\cite{dong2024flex}. Further details of model architectures and hyperparameter selection are provided in Appendix~\ref{sec:training}.

To address the varying training token budgets induced by different linearizations, we compare SLMs with pause-token baselines~\cite{goyal2024think}: $\textsc{Pause-}K = (x_1, p_1, \ldots, p_K, x_2, p_1, \ldots, p_K, x_3, \ldots)$, where $p_k$ are learnable embedding vectors. For $K=1$, the computational load is roughly comparable to $\textsc{Lin}^1$ ($2.4\times$); for $K=2$, it is roughly comparable to $\textsc{Lin}^2$ ($3.2\times$).

\subsection{Evaluation Methods}

\paragraph{Document-Level Language Modeling.}
\label{sec:doc-ppl}
Our SLMs model a joint distribution $p(\mathbf{x}, \mathbf{y})$ over sentences and parse trees (Eq.~\ref{eq:joint}). The performance of language modeling is measured by the marginal probability of the sentence: $p(\mathbf{x}) = \sum_{\mathbf{y}} p(\mathbf{x}, \mathbf{y})$.
Since summing over all possible parse trees is intractable, we follow \citet{zhao-etal-2025-systematic} and approximate the marginal using a parser as the proposal distribution:
\begin{equation}
p(\mathbf{x}) \;\approx\; \sum_{k=1}^{K} p(\mathbf{x}, \mathbf{y}_k),
\qquad \mathbf{y}_k \sim q(\mathbf{y} \mid \mathbf{x})
\label{eq:marginal}
\end{equation}
\noindent where $q(\mathbf{y} \mid \mathbf{x})$ is the constituency parser and $K=300$ (the top-300 parse trees by parser score).
For each sampled tree $\mathbf{y}_k$, we linearize it into $\mathbf{z}_k = \textsc{Linearize}(\mathbf{x}, \mathbf{y}_k)$ and compute $p(\mathbf{x}, \mathbf{y}_k) = \prod_{t=1}^{T_k} p(z_{k,t} \mid z_{k,<t})$ autoregressively.
Let $\mathbf{X} = (\mathbf{x_1}, ..., \mathbf{x_N})$ be a document $\mathbf{X}$ with $N$ sentences. 
Document-level perplexity for a SLM is then:
\begin{equation}
\text{PPL}_{\text{syn}}(\mathbf{X}) = \exp\!\left( -\frac{1}{N_w} \sum_{i=1}^{N} \log p(\mathbf{x}_i \mid \mathbf{x}_{<i}) \right)
\label{eq:doc-ppl}
\end{equation}
\noindent where $p(\mathbf{x}_{i} | \mathbf{x}_{< i}) = p(\mathbf{x}_{\le i}) / p(\mathbf{x}_{< i})$ is computed via Eq.~\ref{eq:marginal} and $N_w$ is the number of terminals in the document. 

In practice, computing this conditional probability would require marginalizing over all $i$ parse trees (each with 300 samples), which is computationally prohibitive. 
Therefore, following \citet{sartran-etal-2022-transformer}, we approximate it by greedy selection: for each of the preceding $i-1$ sentences, we choose the single tree $\mathbf{y}$ that maximizes $p(\mathbf{x}, \mathbf{y})$ and use it as a single-path prefix for the $i$-th sentence.

This tree-marginalized perplexity compares SLMs and terminal baselines on equal footing. %for $\textsc{Lin}^0$ (terminal), there is no tree structure to marginalize over ($\mathbf{y}$ is vacuous) and Eq.~\ref{eq:doc-ppl} reduces to standard autoregressive perplexity. 
%For pause-token models ($\textsc{Lin}^{\text{pause}}_K$), there is similarly no parse tree---the model defines $p(\mathbf{x})$ directly over the flat sequence $\mathbf{z} = (\mathbf{x}, p_1, \ldots, p_K)$, and Eq.~\ref{eq:doc-ppl} reduces to standard perplexity computed over terminal tokens.

\paragraph{Syntactic evaluations.} We probe syntactic competence with two benchmarks. SG~\cite{hu-etal-2020-systematic} tests six targeted generalization subtasks, for which we decode with the word-synchronous beam search method of \citet{stern-etal-2017-effective}; BLiMP~\cite{warstadt-etal-2020-blimp-benchmark} evaluates whether models assign higher likelihood to grammatical sentences in minimal pairs across 12 grammatical phenomena.

\paragraph{Downstream evaluations.} We measure downstream performance with task-specific fine-tuning on XSum summarization~\cite{narayan-etal-2018-dont}, reporting ROUGE-1/2/L~\cite{lin-2004-rouge}, and on BoolQ~\cite{clark-etal-2019-boolq} for boolean question answering. For FineWeb-Edu models, we follow the OLMES protocol~\cite{gu-etal-2025-olmes} and evaluate an 11-task suite spanning WinoGrande~\cite{sakaguchi2020winogrande}, HellaSwag~\cite{zellers-etal-2019-hellaswag}, BoolQ, MMLU~\cite{hendrycks2021measuring}, MMLU-Redux~\cite{gema-etal-2025-done}, OpenBookQA~\cite{mihaylov-etal-2018-suit}, CommonsenseQA~\cite{talmor-etal-2019-commonsenseqa}, SocialIQA~\cite{sap-etal-2019-social}, ARC-Easy and ARC-Challenge~\cite{clark2018think}, and PIQA~\cite{bisk2020piqa}.

\begin{table*}[htbp]
  \centering
  \setlength{\tabcolsep}{9pt}
  \begin{tabular}{lccccc}
    \toprule
    Model & XSum R-AVG$\uparrow$ & BoolQ$\uparrow$ & SG$\uparrow$ & BLiMP$\uparrow$ & Doc-PPL$^\ddagger$$\downarrow$ \\
    \midrule
    Terminal-100M & 21.99 & 67.40 & 69.98 & 70.77 & 10.19 \\
    Tree-100M & \textbf{22.82} & 67.86 & 81.64 & 80.90 & 10.17 \\
    TGTree-100M & 22.69 & 68.38 & 79.40 & 81.48 & 10.36 \\
    TG-100M & 20.13 & 64.31 & 78.80 & 83.23 & 11.20 \\
    TGNomask-100M & 21.60 & 67.25 & 78.47 & 81.77 & 10.14 \\
    TGNomask-Aug-100M & 22.04 & 68.81 & 78.44 & \textbf{83.64} & 10.24 \\
    Tree-NoONT-100M & 22.29 & 68.13 & \textbf{85.37} & 81.00 & \textbf{10.08} \\
    Tree-Compress-100M & 22.63 & 68.04 & 78.18 & 81.44 & 10.09 \\
    Tree-TripleCNT-100M & 21.18 & \textbf{69.02} & 79.83 & 83.38 & 10.41 \\
    Tree-Shuffle-100M & 21.76 & 67.77 & 76.96 & 67.67 & 13.69 \\
    TGNomask-Mix-TG-100M & 21.86 & 67.49 & 75.96 & 83.15 & 10.16 \\
    TGTree-Mix-TG-100M & 22.25 & 67.74 & 79.25 & 82.76 & 10.16 \\
    Pause-1-100M & 19.86 & 66.06 & 78.47 & 69.99 & 10.64 \\
    \midrule
    Terminal-500M & 20.71 & 69.97 & \textbf{71.25} & 64.74 & \textbf{3.06} \\
    Tree-500M & \textbf{22.05} & \textbf{71.90} & 63.26 & 76.07 & 3.32\\
    TGTree-500M & 21.96 & 70.86 & 66.50 & \textbf{77.63} & 3.40\\
    TGNomask-Aug-500M & 21.63 & 71.04 & 57.73 & 75.99 & 3.32\\
    \bottomrule
  \end{tabular}
  \caption{Key 100M and 500M results on BBC News. Best per metric in \textbf{bold}. $\ddagger$ Tree linearization model PPLs are upper bounds.}
  \label{tab:exp1-results}
\end{table*}

\subsection{Experiment 1: Disentangling Tree Linearization and Attention}
\label{sec:exp1}

For training corpus construction, we use the Benepar parser~\cite{kitaev-klein-2018-constituency}, which is trained on news-domain data, to obtain constituency parses for BBC News articles filtered from each subset of the FineWeb dataset~\cite{penedo2024fineweb}. 
The resulting corpus covers approximately 300M articles from BBC News domains and contains $\sim$10B terminal tokens.
We then train all 100M-parameter variants on this corpus for one epoch and evaluate them on XSum ROUGE, BoolQ, SG, and BLiMP.
Our goal is to identify which tree linearization and which attention mechanisms drive the syntactic benefit, and to screen variants for scaling.
Table~\ref{tab:exp1-results} reports results for all representative variants.

\paragraph{Structual Attention vs.\ Causal Attention.}
SLMs with causal attention outperform structural attention models on downstream tasks.
On XSum, Tree and TGTree (22.82, 22.69) exceed TGNomask-Aug (22.04), TG (20.13), and the Terminal baseline (21.99); BoolQ shows the same ordering.
The full STACK/COMPOSE mask (TG) underperforms even the terminal baseline on both tasks, while a softer variant (TGNomask-Aug) yields a partial recovery.
Structural attention models exhibit a clear syntax--utility tradeoff: TG and its variants lead only on BLiMP (up to 83.64), where the hard constraints of STACK/COMPOSE attention enforce grammatical consistency, but the same rigidity penalizes general language modeling.
Causal attention models suffer no such tradeoff---Tree achieves the best XSum and the second-best SG (81.64) while TGNomask-Aug manages only 78.44.
The mask is neither necessary nor beneficial for most objectives. %; the tree input format alone suffices.

\paragraph{Hybrid Attention Fails to Improve.}
Hybrid attention variants (TGNomask-Mix-TG, TGTree-Mix-TG) yield no benefit over causal Tree SLM, adding complexity without improving performance.

\paragraph{Document-Level Perplexity.}
%The tree-marginalized document-level PPL (Eq.~\ref{eq:doc-ppl}) provides an orthogonal signal that tests pure language modeling quality independent of task-specific decoding.
%At 100M (full results in Appendix~\ref{sec:full-100m}), 
The PPL of Tree matches Terminal within 0.02 nats, confirming that linearizing trees into the input sequence does not degrade the model's core LM capability.
The STACK/COMPOSE mask is the exception: TG degrades PPL to 11.20, consistent with its downstream underperformance.
TGNomask is competitive with causal tree models, supporting \citet{zhao-etal-2025-systematic}'s finding that removing sub-constituent masks improves LM quality.
Among all variants, Tree-NoONT achieves the lowest PPL, suggesting ONTs impose a small but detectable LM cost, while Tree-Shuffle collapses to 13.69, confirming that scrambled labels and non-structural information destroy both LM and downstream performance.
Pause-1 underperform the terminal baseline on document-level PPL at this scale.
Overall, PPL rankings align with downstream task rankings, reinforcing the conclusion that tree linearization with causal attention is the best choice.

% \paragraph{Tree Linearization Ablations} Within all SLMs with causal attention, Opening non-terminals primarily aid generation (higher XSum with ONTs), not syntactic discrimination (SG is higher without them), suggesting they help the model anticipate upcoming constituent structure.
% Random NT labels (Tree-Shuffle) cause a catastrophic BLiMP collapse, confirming that the model exploits syntactic label identity rather than only tree topology.
% Doubling CNTs (TGTree, 2$\times$) provides the best syntax--generation balance; tripling (3$\times$) further improves BLiMP at the cost of degraded generation.
% Merging consecutive identical CNTs (Tree-Compress) reduces SG, indicating that each constituent's closure independently contributes syntactic signal.

\paragraph{Tree Linearization Ablations. }
The experiments on all SLMs with causal attention form an ablation study on tree linearization.
Tree-NoONT improve syntactic discrimination but impair generation, suggesting ONTs primarily help the model anticipate upcoming constituent structures rather than enforce syntactic constraints. Randomizing tree linearization in Tree‑Shuffle causes a catastrophic collapse on BLiMP, confirming that SLMs exploit syntactic tree topology and label identity. Doubling CNTs (TGTree) achieves the best trade‑off between syntax and generation; tripling CNTs (Tree-TripleCNT) further improves BLiMP at the cost of generation quality. Merging consecutive identical closing non-terminals as in Tree‑Compress reduces syntactic generalization, indicating that each constituent closure contributes independently to the syntactic signal. 

\paragraph{Scaling Up.} After the 100M experiments, we select four variants and train them at 500M parameters on BBC News to verify that the advantage of tree linearization with causal attention persists before committing to larger-scale experiments.
Table~\ref{tab:exp1-results} shows the results, confirming that Tree and TGTree maintain their downstream lead over the terminal baseline and the structural attention variant.
SG scores decline for all models at this scale, which we discuss in \S\ref{sec:sg-discussion}.
% We proceed to 1B with Tree and TGTree as the primary syntactic variants.

There is a practical note on throughput: SLMs with causal attention benefit from FlashAttention compatibility, whereas SLMs with structural attention require slower custom attention implementations.
Table~\ref{tab:efficiency} reports measured training throughput.

\begin{table}[tbp]
  \centering
  \fontsize{10pt}{12pt}\selectfont
  \setlength{\tabcolsep}{6pt}
  \begin{tabular}{lccc}
    \toprule
    Scale & Causal Mask & Structural Mask & Ratio \\
    \midrule
    100M & 62,175 token/s & 37,975 token/s & 1.6$\times$ \\
    500M & 12,550 token/s & 8,316 token/s & 1.5$\times$ \\
    \bottomrule
  \end{tabular}
  \caption{Training throughput on H800 GPUs. Causal models use FlashAttention-2; structural mask models use FlexAttention. Mixing-head mask: 3,375 token/s (100M scale, 18$\times$ slower than causal).}
  \label{tab:efficiency}
\end{table}

The gap persists with scale, and mixing-head mask models are 18$\times$ slower.
Combined with the performance evidence above, the causal tree approach is the clear practical choice for scaling.

Based on these results, we select tree linearization Tree ($\textsc{Lin}^1$) and TGTree ($\textsc{Lin}^2$) , both with causal attention, as the primary variants for scaling, dropping all mixing-head and structural attention approaches.

\subsection{Experiment 2: Scaling to 1B on FineWeb-Edu --- Main Result}

FineWeb-Edu-100BT~\cite{penedo2024fineweb} is a $\sim$100B-token sampled subset of the FineWeb-Edu corpus, filtered for educational quality and spanning diverse domains.
We train five models at 1B parameters on this data: Terminal, Tree, TGTree, Pause-1 and Pause-2.
Evaluation uses the 11-task OLMES suite~\cite{gu-etal-2025-olmes} with both native-format and terminal-format scoring (\S\ref{sec:terminal-format-evaluation}). Table~\ref{tab:main-result} reports per-task accuracy for all model settings.

\begin{table*}[htbp]
  \centering
  %\small
  \setlength{\tabcolsep}{1.5pt}
  \begin{tabular}{lccccccccccccc}
    \toprule
    Model & WG$\uparrow$ & HS$\uparrow$ & BQ$\uparrow$ & MM$\uparrow$ & MMR$\uparrow$ & OBQA$\uparrow$ & CSQA$\uparrow$ & SIQA$\uparrow$ & ARCE$\uparrow$ & ARCC$\uparrow$ & PIQA$\uparrow$ & \textbf{AVG}$\uparrow$ \\
    \midrule
    Terminal & 58.25 & 61.46 & 59.11 & 37.01 & 37.39 & 43.60 & 53.15 & 49.13 & 69.47 & 43.14 & 73.83 & 53.23 \\
    Tree & 62.35 & 59.30 & 64.83 & 36.91 & 37.32 & 41.20 & 56.84 & 47.24 & 67.02 & 46.15 & 74.10 & 53.93 \\
    TGTree & \textbf{64.17} & 58.35 & \textbf{68.20} & 37.95 & 38.23 & 38.20 & \textbf{58.56} & 46.47 & 67.02 & 44.15 & 73.56 & 54.08 \\
    Tree$_{\text{term}}$ & 63.46 & 62.76 & 64.83 & 37.28 & 37.49 & 43.80 & 56.27 & 48.26 & 69.12 & 44.48 & 74.16 & 54.72 \\
    TGTree$_{\text{term}}$ & 63.14 & 63.37 & \textbf{68.20} & 38.23 & 38.49 & \textbf{45.80} & 58.15 & 50.31 & 69.65 & 44.82 & 74.92 & \textbf{55.92} \\
    Pause-1 & 58.25 & 63.18 & 63.91 & 38.11 & 38.60 & 43.00 & 55.94 & \textbf{50.97} & \textbf{73.68} & 43.81 & \textbf{75.03} & 54.95 \\
    Pause-2 & 63.30 & \textbf{64.62} & 61.62 & \textbf{38.46} & \textbf{39.26} & 44.60 & 56.84 & 48.98 & 72.11 & \textbf{47.83} & 74.86 & 55.68 \\
    \bottomrule
  \end{tabular}
  \caption{Per-task accuracy at 1B FineWeb-Edu. WG = WinoGrande, HS = HellaSwag, BQ = BoolQ, MM = MMLU, MMR = MMLU-Redux, OBQA = OpenBookQA, CSQA = CommonsenseQA, SIQA = SocialIQA, ARCE = ARC-Easy, ARCC = ARC-Challenge. Best per task in \textbf{bold}. }
  \label{tab:main-result}
\end{table*}

\paragraph{Main result.}
TGTree$_{\text{term}}$ achieves the highest 11-task average (55.92, +2.69 over Terminal).
Tree and TGTree exceed the Terminal baseline, confirming that the syntactic tree linearization advantage persists at 1B scale on diverse corpora and on diverse downstream tasks.
The 2.69-point gap on FineWeb-Edu is substantially larger than any BBC-scale gap, supporting the hypothesis that explicit structural modeling provides greater benefit as data diversity increases.

\paragraph{Terminal-format evaluation is essential.}
Tree$_{\text{term}}$ (54.72) exceeds Tree (53.93) by 0.79 points, and TGTree$_{\text{term}}$ (55.92) exceeds TGTree (54.08) by 1.84 points.
These gaps are purely an evaluation artifact: without terminal-format scoring, the true contribution of tree structure is substantially understated.
We recommend terminal-format evaluation as a standard protocol for future SLM research.

\paragraph{Explicit Structure vs.\  Dummy Tokens.}
TGTree$_{\text{term}}$ slightly outperforms Pause-2 by 0.24 points overall, which we interpret as broadly competitive performance with a slight edge.
The task-level pattern reveals complementarity rather than dominance: Pause-2 leads on ARC-Challenge (+4.69 over Terminal) and HellaSwag (+3.16), while TGTree$_{\text{term}}$ leads on BoolQ (+9.09) and CommonsenseQA (+5.00). The largest gains for TGTree$_{\text{term}}$ occur on three tasks requiring multi-hop integration---BoolQ (+9.09), CommonsenseQA (+5.00), and WinoGrande (+4.89)---precisely where constituency structures may help organize context across spans. 
Explicit syntactic structures and learned latent computation appear to benefit different reasoning types; a combination of the two is a natural direction for future work. 

\subsection{Scaling Trend Summary}

\begin{table*}[htbp]
  \centering
  \setlength{\tabcolsep}{12pt}
  \begin{tabular}{lllcc}
    \toprule
    Scale & Data & Metric & Terminal & Best SLM ($\Delta$) \\
    \midrule
    100M & BBC & XSum R-AVG & 21.99 & Tree: 22.82 (+0.83) \\
    500M & BBC & XSum R-AVG & 20.71 & Tree: 22.05 (+1.34) \\
    1B & BBC & XSum R-AVG & 20.31 & Tree: 21.30 (+0.99) \\
    1B & FW-Edu & 11-task AVG & 53.23 & TGTree$_{\text{term}}$: 55.92 (+2.69) \\
    \bottomrule
  \end{tabular}
  \caption{Scaling trend across model sizes and datasets. $\Delta$ indicates the performance gain of the best SLM over the Terminal baseline (Best SLM --- Terminal).}
  \label{tab:scaling-summary}
\end{table*}

Table~\ref{tab:scaling-summary} confirms that the advantage beyond the baseline Terminal is consistently positive across all scales and datasets.
On BBC, the gap between Terminal and best SLM is non-monotonic (+0.83 $\rightarrow$ +1.34 $\rightarrow$ +0.99; the 1B BBC run was a supplementary XSum-only evaluation), possibly reflecting saturation on the narrow 10B-token corpus.
On diverse FineWeb-Edu data, the gap widens to +2.69 on 11-task AVG, consistent with the hypothesis that structures help more when the data distribution is more varied.

% \begin{table}[htbp]
%   \centering
%   \resizebox{\columnwidth}{!}{%
%   \begin{tabular}{lccc}
%     \toprule
%     Variant & XSum$\uparrow$ & SG$\uparrow$ & BLiMP$\uparrow$ \\
%     \midrule
%     Tree (baseline) & \textbf{22.82} & 81.64 & 80.90  \\
%     Tree-NoONT & 22.29 & \textbf{85.37} & 81.00  \\
%     Tree-Compress & 22.63 & 78.18 & 81.44  \\
%     Tree-Shuffle & 21.76 & 76.96 & 67.67  \\
%     TGTree ($2\times$CNT) & 22.69 & 79.40 & 81.48 \\
%     Tree-TripleCNT ($3\times$CNT) & 21.18 & 79.83 & \textbf{83.38}\\
%     \bottomrule
%   \end{tabular}%
%   }
%   \caption{Ablation results at 100M BBC News.}
%   \label{tab:ablations}
% \end{table}

\section{Discussion}

\subsection{Why Does Tree Linearization Help Without Structural Attention?}
\label{sec:why-works}

We propose three complementary hypotheses:

\paragraph{Implicit multi-task learning.} The model jointly predicts terminal tokens and syntactic structures, making constituency boundaries explicit through non-terminals rather than leaving them to be inferred from co-occurrence statistics. \citet{sartran-etal-2022-transformer} attribute the perplexity gains of such joint modeling to the ``additional syntactic supervision'' that non-terminals provide by constraining admissible word classes at each position.

\paragraph{Learned compositional bottleneck.} CNT positions create natural ``summary points'' where the model must compress constituent information. Unlike TG's hard-coded COMPOSE attention, this bottleneck is learned and flexible---the model chooses what information to route through CNT positions based on linguistic context. This flexibility becomes more valuable with more diverse training data, which may explain the larger FineWeb-Edu advantage, suggesting that a soft bias offers more than a hard constraint.

\paragraph{Adaptive computation allocation.} Pause-$K$ controls for extra tokens by allocating a fixed number of learned latent tokens after every terminal. Tree models instead use parse-dependent structural tokens: richer constituency structures introduce more ONT/CNT positions at phrase boundaries. This adaptive allocation may help terminal prediction by placing extra computation where composition occurs. The fact that TGTree$_{\text{term}}$ is competitive with and modestly exceeds Pause-2 on FineWeb-Edu suggests that tree structures may supply adaptive computation.

%\paragraph{Soft bias $>$ hard constraint.} Causal attention lets the model learn \textit{when} to use syntactic structure; STACK/COMPOSE forces it even when harmful. The finding that TG (hard mask) underperforms on general tasks while excelling at BLiMP supports this: the hard mask over-constrains attention, trading general capability for grammatical sensitivity. Causal tree models get the syntactic benefit without the rigidity penalty.

\subsection{The SG Degradation Puzzle}
\label{sec:sg-discussion}

SG for tree models drops from 81.64 (100M) to 63.26 (500M) on BBC.
We consider three non-exclusive explanations: (1) \textbf{Domain overfitting}: BBC News is narrow ($\sim$10B tokens); tree models see 2.4$\times$ more tokens per raw document, potentially overfitting to the domain's syntactic patterns.
(2) \textbf{Benchmark limitations}: SG has known ceiling effects (Center Embedding: 87.5--94.6 across all 100M models) and may not reliably discriminate above certain capability thresholds.
(3) \textbf{Capability integration}: At larger scales, models may integrate syntax with semantics rather than treating it as an isolated phenomenon, and may trade pure syntactic discrimination for integrated linguistic competence.

We treat downstream task performance as the more reliable signal at scale.

\section{Conclusion}

We scale up the empirical study of syntactic language models and systematically disentangle tree linearization from attention mechanism. Across our 12 SLM variants, with model sizes up to 1B parameters, training data up to 100B tokens, and diverse downstream and syntactic evaluations, we find that the main benefit comes from tree linearization rather than hard structural masks. Our experiments at larger scales further show that explicit syntactic structure remains useful beyond small or narrowly syntactic settings. At the same time, SLMs with standard causal attention preserve FlashAttention compatibility and introduce no architectural overhead. These results support a practical pipeline for syntax-augmented pretraining: parse text, linearize constituency structure, and train standard Transformers over the resulting sequences.



\section*{Limitations}

Both training-data parsing and input preprocessing require a parser to construct linearized tree inputs. This adds preprocessing cost, and parser errors may propagate into the training signal. We do not quantify the impact of parse quality. In addition, we investigate only English, leaving syntactic structures in other languages unexplored.

Our discussion of learned CNT bottlenecks and adaptive computation is explanatory rather than causal. The core conclusions are supported by controlled comparisons across linearizations, attention masks, pause-token baselines, and scales, but we do not directly intervene on internal representations or attention dynamics to prove how tree tokens are used. Such mechanistic probes are left to future work.

% Acknowledgments (anonymized)
\section*{Acknowledgments}
This work was supported by the HPC Platform of ShanghaiTech University, the SIST Computing Platform and the Core Facility Platform of Computer Science and Communication, SIST, ShanghaiTech University. 


\bibliography{anthology-1,anthology-2,custom}

\appendix

\section{Training Details}
\label{sec:training}

% TODO: 
\subsection{Implementation and Hyperparameters}

We use an OLMo-based~\cite{groeneveld-etal-2024-olmo} decoder-only Transformer with SwiGLU~\cite{shazeer2020gluSWIGLU} activation, RoPE~\cite{su2024roformerROPE} positional encoding, RMSNorm~\cite{zhang2019rootRMS} for pre-normalization, and weight tying.
Under this architecture, we train different SLMs at 100M, 500M, and 1B parameters. Table~\ref{tab:scales} lists the three scale configurations.

\begin{table}[htbp]
  \centering
  \begin{tabular}{lcccc}
    \toprule
    Scale & $d_\text{model}$ & Layers & Heads & Params \\
    \midrule
    100M & 768 & 12 & 12 & $\sim$100M \\
    500M & 1408 & 16 & 16 & $\sim$500M \\
    1B & 2048 & 16 & 16 & $\sim$1B \\
    \bottomrule
  \end{tabular}
  \caption{Model configurations at each scale.}
  \label{tab:scales}
\end{table}

We train all models from scratch for one epoch. Training uses the AdamW optimizer~\cite{loshchilov2017decoupledADAM} with $\beta_1{=}0.9$, $\beta_2{=}0.95$, $\epsilon{=}10^{-8}$, weight decay 0.1, and gradient clipping at norm 1.0. We use a cosine learning rate scheduler with 2,000 warmup steps; after the learning rate decays to $10^{-5}$, it is held at that value for the rest of training. The maximum sequence length is 2,048 tokens. The learning rate $\eta$ and batch size $B$ are set according to Step Law~\cite{li2025predictable}:

\begin{equation}
    \begin{gathered}
       \eta(N,D) = 1.79N^{-0.713}D^{0.307} \\
B(D) = 0.58D^{0.571}
    \end{gathered}
\label{eq:steplaw}
\end{equation}
where $N$ is the number of non-embedding model parameters and $D$ is the size of training dataset. 

\subsection{Computational Costs}
The 100M-parameter models were trained on 4 H800 GPUs for 133 hours. The 500M-parameter models were trained on 4 H800 GPUs for 125 hours. The 1B-parameter models were trained on 64 A800 GPUs for 385 hours (approximately 16.0 days).

\subsection{Tokenizer and Structural Tokens}
We use tokenizers with reserved non-terminal tokens (e.g., \texttt{NP)}, \texttt{(S}, \texttt{VP)}).
The tokenizer is based on the GPT-2 tokenizer~\cite{radford2019language} or the Qwen3 tokenizer~\cite{yang2025qwen3}.
For tree linearization models, opening non-terminals (ONT) and closing non-terminals (CNT) are represented as reserved tokens so that linearized constituency structure is preserved during pretraining and evaluation.

\section{Evaluation Setup Details}

\subsection{Downstream Fine-Tuning on BBC News: XSum and BoolQ}
On BBC News, both XSum and BoolQ are evaluated via task-specific fine-tuning from the pretrained checkpoint (optimizer state reset).
All models are fine-tuned with AdamW ($\beta_1{=}0.9$, $\beta_2{=}0.95$, weight decay 0.1) and a cosine learning rate schedule with 100 warmup steps decaying to $10^{-6}$.

\paragraph{XSum.} Fine-tuned for 3 epochs with peak LR $5{-}6\times10^{-5}$ (depending on variant), global batch size 40.
During evaluation, summaries are generated via beam search (beam size 6, max 150 tokens).
For tree linearization models, beam search proceeds at the word level (word-sync) with a maximum of 75 word steps and a grammar constraint of at most 4 consecutive opening non-terminals.
ROUGE-1, ROUGE-2, and ROUGE-L are computed with the standard evaluate library; R-AVG denotes their arithmetic mean.

\paragraph{BoolQ.} Fine-tuned for 5 epochs with peak LR $3\times10^{-4}$, global batch size 40.
Accuracy is measured on the validation split.

\subsection{Syntactic Generalization (SG)}
SG evaluation follows Zhao et al.~\cite{zhao-etal-2025-systematic}.
For terminal and pause-token models, sentence-level log-probability is computed via a direct forward pass summing the cross-entropy loss over the target span.
For tree linearization models, we use word-sync beam search (beam size 300) with a DFS expansion strategy: at each word step, non-terminal tokens are expanded until the beam stabilizes or a maximum of $1.2\times$ the number of generated non-terminal tokens is reached, after which terminal tokens advance the beam.
Conditional probability scores are compared using the task-specific formulas from Hu et al.~\cite{hu-etal-2020-systematic} to produce per-task accuracy.

\subsection{BLiMP}
For each BLiMP minimal pair, we sample $K=300$ parse trees from the Benepar parser~\cite{kitaev-klein-2018-constituency} and compute the tree-marginalized probability of each sentence via $p(\mathbf{x}) \approx \sum_{k=1}^{K} p(\mathbf{x}, \mathbf{y}_k)$, following the same marginalization procedure as document-level PPL (Eq.~\ref{eq:marginal}).
A pair is scored correct if the good sentence receives higher marginalized probability than the bad sentence.
Each of the 67 subtasks contains 1,000 pairs.
For terminal and pause-token models, $K=1$ (no tree marginalization).

\subsection{OLMES Multiple-Choice Evaluation on FineWeb-Edu}
All 11 OLMES tasks (HellaSwag, WinoGrande, BoolQ, ARC-Easy, ARC-Challenge, PIQA, MMLU, MMLU-Redux, OpenBookQA, SocialIQA, CommonsenseQA) are evaluated following the OLMES completion formulation standard~\cite{gu-etal-2025-olmes}. 
As in OLMES, each task uses a few-shot setting: HellaSwag (5-shot), WinoGrande (5-shot), BoolQ (3-shot), ARC-Easy (3-shot), ARC-Challenge (3-shot), PIQA (3-shot), MMLU (5-shot), MMLU-Redux (5-shot), OpenBookQA (5-shot), SocialIQA (3-shot), CommonsenseQA (3-shot).
The model scores each answer option by summing log-probability over the continuation; for terminal-format evaluation, only terminal tokens contribute to the score (Eq.~\ref{eq:decom}).

\section{SG Subtask Breakdown}

Table~\ref{tab:sg-subtasks} breaks down SG performance by subtask for key 100M variants.

\begin{table}[htbp]
  \centering
  \fontsize{8pt}{10pt}\selectfont
  \setlength{\tabcolsep}{1.8pt}
  \begin{tabular}{lcccccc}
    \toprule
    Model & Agr$\uparrow$ & CtrEmb$\uparrow$ & GrdnPath$\uparrow$ & GrsSyn$\uparrow$ & Lic$\uparrow$ & LDD$\uparrow$ \\
    \midrule
    Terminal & 39.66 & 87.50 & 78.29 & 84.78 & 61.28 & 68.35 \\
    Tree & 81.03 & 89.29 & 84.21 & 83.70 & 80.08 & 71.56 \\
    TGTree & 67.24 & 94.64 & 82.24 & 88.04 & 69.92 & 74.31 \\
    TG & 79.31 & 92.86 & 81.58 & 79.35 & 71.80 & 67.89 \\
    Tree-NoONT & 87.93 & 89.29 & 85.53 & 97.83 & 80.08 & 71.56 \\
    \bottomrule
  \end{tabular}
  \caption{SG subtask breakdown at 100M. Agr = Agreement, CtrEmb = Center Embedding, GrdnPath = Garden Path Effects, GrsSyn = Gross Syntactic Expectation, Lic = Licensing, LDD = Long Distance Dependencies. $\uparrow$ = higher is better.}
  \label{tab:sg-subtasks}
\end{table}

Agreement shows the largest improvement (Tree +41.4 vs.\ Terminal). Center Embedding shows ceiling effects across all models (87.5--94.6). Long-Distance Dependencies show the smallest gains (3--6 points).


\section{Non-Terminal Token Score Decomposition}
\label{sec:decom}

We analyze how non-terminal (ONT/CNT) tokens affect multiple-choice evaluation by decomposing the MC score into terminal and non-terminal components (Eq.~\ref{eq:decom} in \S\ref{sec:terminal-format-evaluation}).
Table~\ref{tab:decom} reports per-task results for syntactic models at 1B on FineWeb-Edu (Qwen3 tokenizer).

\begin{table*}[htbp]
  \centering
  \begin{tabular}{lccc}
    \toprule
    Task & $\Delta$Acc (Term $-$ Full) & Flip Rate & NT log-p Gap \\
    \midrule
    HellaSwag & +3.36 & 20.1\% & +2.23 \\
    WinoGrande & +0.55 & 8.5\% & +2.18 \\
    ARC-Easy & +2.46 & 17.2\% & +3.24 \\
    ARC-Challenge & $-$1.67 & 22.1\% & +2.89 \\
    PIQA & +0.22 & 11.2\% & +2.46 \\
    SocialIQA & +0.71 & 24.9\% & +3.57 \\
    CommonsenseQA & $-$1.14 & 22.4\% & +5.98 \\
    OpenBookQA & +1.80 & 21.6\% & +3.77 \\
    \midrule
    Mean & +0.79 & 18.5\% & +3.29 \\
    \bottomrule
  \end{tabular}
  \caption{Decomposed MC accuracy for Tree-1B (FineWeb-Edu). $\Delta$Acc = terminal $-$ full accuracy; Flip Rate = percentage of examples whose answer changes between full and terminal scoring; NT log-p gap = mean non-terminal log-probability minus mean terminal log-probability. Tasks without NTs (BoolQ, MMLU, MMLU-Redux) excluded.}
  \label{tab:decom}
\end{table*}

Three findings emerge.
First, syntactic models assign substantially higher probability to non-terminal tokens than to terminal tokens---the NT log-probability gap averages +3.3 nats ($\sim$27$\times$ higher per-token probability for the Qwen3 tokenizer; the GPT-2 tokenizer yields a +2.7 nat gap, $\sim$15$\times$).
This means the full MC score is dominated by bracket prediction quality rather than content prediction.

Second, removing NTs from the score changes 18.5\% of MC answers on average (range 8.5--24.9\% across tasks).
For the terminal baseline, flip rates are near-zero ($\le$0.4\%), confirming that the decomposition correctly isolates the effect of non-terminals.

Third, flips are roughly symmetric: across tasks, 7.2\% flip to correct (NT removal helps) and 6.4\% flip to wrong (NT removal hurts).
This symmetry suggests that non-terminals introduce noise rather than systematic bias---they help and hurt in roughly equal measure.
Nevertheless, terminal-format evaluation is methodologically necessary to ensure fair comparison: without it, syntactic models are penalized on tasks where their intrinsic bracket structure preference conflicts with the MC scoring objective.

\end{document}